%%%%%%%%%%  scan idx 582  |  folio 569  %%%%%%%%%%
\SGAPrintPage{569}%

Now these two squares are the ``base-change'' squares obtained from the two squares (7.10.0), so that one can apply (IV 3). One need only verify that these squares correspond to pairs \((X,Y_t)\) and \((X,Y_s)\) of schemes over \(Y\) that are Tor-independent over \(Y\). This is trivial for the first pair, since \(Y_t\to Y\) is flat, and also immediate for the second, since the immersion \(Y_s\to Y\) is regular of codimension 1 and remains so after the base change \(X\to Y\).

Now no longer suppose that \(g\) is proper, but suppose instead that it has finite Tor-dimension. I claim that the same is then true of \(g_t\) and \(g_s\). This is trivial for the first, and for the second follows from the fact (just noted) that \(X\) and \(Y_s\) are Tor-independent over \(Y\). One can therefore consider the square of homomorphisms
\sgaStableTarget{sga6.e10.eq.7-10-2}\begin{equation}\tag{7.10.2}
\begin{tikzcd}[column sep=large,row sep=large]
K_\bullet(Y_t) \arrow[r,"g_t^*"] \arrow[d,"\sigma_Y"'] &
K_\bullet(X_t) \arrow[d,"\sigma_X"]\\
K_\bullet(Y_s) \arrow[r,"g_s^*"] &
K_\bullet(X_s),
\end{tikzcd}
\end{equation}
where the horizontal arrows are defined in (IV 2.12). Note that, when \(X\) and \(Y\) satisfy the conditions of 7.7, i.e., have regular fibers, \(g\) automatically has finite Tor-dimension (since \(Y\) is regular), and the preceding diagram can be interpreted as a commutative diagram of rings \(K^\bullet\).

To prove the commutativity of (7.10.2), it again suffices to prove that of the front squares of the prism analogous to \((*)\), but with the direct-image homomorphisms \(g_{t*}\), \(g_*\), \(g_{s*}\) replaced by the inverse-image homomorphisms \(g_t^*\), \(g^*\), \(g_s^*\). Now the commutativity of these squares follows at 
%%%%%%%%%%  scan idx 583  |  folio 570  %%%%%%%%%%
\SGAPrintPage{570}%
once from that of the corresponding squares (7.10.0), taking account of the transitivity of inverse images (IV 2.12).

\sgaStableTarget{sga6.e10.rem.7-10-3}\textbf{Remark 7.10.3.} When \(g\) is proper, the commutativity of (7.10.1) immediately implies that of the following diagram, obtained from it by passage to the associated graded groups:
\sgaStableTarget{sga6.e10.eq.7-10-1bis}\begin{equation}\tag{7.10.1 bis}
\begin{tikzcd}[column sep=large,row sep=large]
\Gr_\bullet(X_t) \arrow[r,"g_{t*}"] \arrow[d,"\sigma_{X\bullet}"'] &
\Gr_\bullet(Y_t) \arrow[d,"\sigma_{Y\bullet}"]\\
\Gr_\bullet(X_s) \arrow[r,"g_{s*}"] &
\Gr_\bullet(Y_s).
\end{tikzcd}
\end{equation}
Likewise, when \(X\) and \(Y\) have regular fibers and are separated, the commutativity of (7.10.2) implies that of
\sgaStableTarget{sga6.e10.eq.7-10-2bis}\begin{equation}\tag{7.10.2 bis}
\begin{tikzcd}[column sep=large,row sep=large]
\Gr^\bullet(Y_t) \arrow[r,"g_t^*"] \arrow[d,"\sigma_Y^\bullet"'] &
\Gr^\bullet(X_t) \arrow[d,"\sigma_X^\bullet"]\\
\Gr^\bullet(Y_s) \arrow[r,"g_s^*"] &
\Gr^\bullet(X_s).
\end{tikzcd}
\end{equation}
When \(X\) and \(Y\) are moreover quasi-projective and smooth over \(S\), one also knows that \(g_t^*\) and \(g_s^*\) are compatible with the filtrations by codimension of supports, and one obtains a commutative diagram like (7.10.2 bis), but with the \(\Gr^\bullet\) replaced by the \(\Gr^\bullet_{\mathrm{top}}\) (cf. Exposé 0, 2.1). One can also note that, when \(X\) and \(Y\) are quasi-projective over \(S\) with regular fibers and separated, \(f_t\) and \(f_s\) are automatically complete-intersection morphisms (VIII 1.1), and one obtains a commutative diagram analogous to (7.10.1 bis), with the \(\Gr_\bullet\) replaced by the \(\Gr^\bullet(\ )_{\mathbb Q}\), taking account of the compatibility result (VIII 3.2).

%%%%%%%%%%  scan idx 584  |  folio 571  %%%%%%%%%%
\SGAPrintPage{571}%

\sgaStableTarget{sga6.e10.para.7-11}\textbf{7.11.} For a scheme \(T\) over an algebraically closed field \(k\), one can define the subgroup of \(K^\bullet(T)\) of elements \emph{algebraically equivalent to zero} as the subgroup of elements of the form
\sgaStableTarget{sga6.e10.eq.7-11-1}\begin{equation}\tag{7.11.1}
  E(x_1)-E(x_0),
\end{equation}
where \(x_0\) and \(x_1\) are two \(k\)-rational points of a proper, regular, connected algebraic curve \(C\) over \(k\), \(E\) is an element of \(K^\bullet(T\times_k C)\), and, for every point \(x\) of \(C\) rational over \(k\), one puts
\[
  E(x)=i_x^*(E)\in K^\bullet(T),\qquad\text{where }
  i_x:T\longrightarrow T\times_k C
\]
is the morphism \(t\mapsto(t,x)\) defined by \(x\). The fact that the elements of the form (7.11.1) form a subgroup of \(K^\bullet(T)\) follows by a standard argument, using a product \(C\times_k C'\) and its pair of points \((x_1,x'_1)\), \((x_0,x'_0)\), and the fact that through two points of an irreducible quasi-projective algebraic scheme over \(k\) (here \(C\times_k C'\)) one can always pass an irreducible algebraic curve. (A reader unwilling to admit this fact may define algebraic equivalence by the group generated by the elements of the form (7.11.1).) The elements algebraically equivalent to zero form an ideal of \(K^\bullet(T)\), as one sees at once, which is moreover stable under the operations \(\lambda^i\) when these are defined, i.e., when one works with the ``naive'' \(K^\bullet\); this ideal is moreover contained in the augmentation ideal of \(K^\bullet(T)\), as one verifies trivially. One can therefore define a quotient ring
\sgaStableTarget{sga6.e10.eq.7-11-2}\begin{equation}\tag{7.11.2}
  K^\bullet_{\mathrm{alg}}(T),
\end{equation}
which is even a \(\lambda\)-ring when one works with naive \(K^\bullet\).

In the same way, one can define a subgroup of elements algebraically equivalent to 0 in \(K_\bullet(T)\), taking \(E\) above in \(K_\bullet(T\times_k C)\), and

%%%%%%%%%%  scan idx 585  |  folio 572  %%%%%%%%%%
\SGAPrintPage{572}%

noting that \(i_x:T\to T\times_k C\) is a \emph{perfect} morphism, since it is obtained by the flat base change \(T\times_k C\to C\) from the morphism \(j_x:x\to C\), which is perfect because \(C\) is regular. One thus obtains a quotient group
\sgaStableTarget{sga6.e10.eq.7-11-3}\begin{equation}\tag{7.11.3}
  K_\bullet^{\mathrm{alg}}(T)
\end{equation}
of \(K_\bullet(T)\). This group is moreover naturally a module over \(K^\bullet_{\mathrm{alg}}(T)\), since one sees at once from the definitions that the product of an element of \(K^\bullet(T)\) with an element of \(K_\bullet(T)\) is algebraically equivalent to zero when one of the two factors is, using (IV 2.12.2) for \(i_x\).

One defines (when working with naive \(K^\bullet\)) the ideal of \(\Gr^\bullet(T)\) of elements algebraically equivalent to zero as the ideal generated by the \(c^i(x)\), where \(i\geq1\) and \(x\in K^\bullet(T)\) is algebraically equivalent to zero; the quotient graded ring
\sgaStableTarget{sga6.e10.eq.7-11-4}\begin{equation}\tag{7.11.4}
  \Gr^\bullet_{\mathrm{alg}}(T)
\end{equation}
is then precisely the Chern ring associated with the \(\lambda\)-ring \(K^\bullet_{\mathrm{alg}}(T)\).

Finally, one defines the graded subgroup of \(\Gr_\bullet(T)\) of elements algebraically equivalent to zero (\(T\) being assumed locally noetherian) as the subgroup whose component of dimension \(i\) is the image in \(\Gr_i(T)\) of the subgroup of \(\Fil_iK_\bullet(T)\) consisting of the elements algebraically equivalent to zero in \(K_\bullet(T)\). One finds at once that the quotient graded group
\sgaStableTarget{sga6.e10.eq.7-11-5}\begin{equation}\tag{7.11.5}
  \Gr_\bullet^{\mathrm{alg}}(T)
\end{equation}
is a module over the ring \(\Gr^\bullet_{\mathrm{alg}}(T)\), by passage to the quotient from the module structure of \(\Gr_\bullet(T)\) over \(\Gr^\bullet(T)\) (3.1).

\sgaStableTarget{sga6.e10.para.7-11-6}\textbf{7.11.6.} The preceding definitions would still make sense without assuming that \(k\) is algebraically closed, taking \(C\) above to be \emph{geometrically} connected over \(k\). It is not clear, however, that these definitions would be very useful; in particular, it does not appear that, in the case of the group \(\Pic(T)\), one

%%%%%%%%%%  scan idx 586  |  folio 573  %%%%%%%%%%
\SGAPrintPage{573}%

recovers the classical ``geometric'' notion of algebraic equivalence considered in 2.4.1. We therefore prefer to say that an element of \(K^\bullet(T)\) is algebraically equivalent to zero when its image in \(K^\bullet(T\otimes_k\bar k)\) is so, where \(\bar k\) is an algebraic closure of \(k\), and to proceed in the same way for the definition of algebraic equivalence in the groups \(K_\bullet(T)\), \(\Gr^\bullet(T)\) and \(\Gr_\bullet(T)\).

\sgaStableTarget{sga6.e10.para.7-12}\textbf{7.12.} With these definitions in place, I claim that \emph{the specialization homomorphism} (7.3.1) \emph{carries elements algebraically equivalent to zero to elements algebraically equivalent to zero}, and consequently defines, by passage to the quotient, a homomorphism
\sgaStableTarget{sga6.e10.eq.7-12-1}\begin{equation}\tag{7.12.1}
  K_\bullet^{\mathrm{alg}}(X_t)\longrightarrow K_\bullet^{\mathrm{alg}}(X_s).
\end{equation}

To see this, one reduces immediately to the case where \(A\) is complete and \(k(s)\) is algebraically closed (using a straightforward compatibility made explicit in 7.15). It suffices to see that \(\sigma_\bullet\) carries an element of the form (7.11.1) of \(K_\bullet(X_t)\) to an element algebraically equivalent to zero in \(K_\bullet(X_s)\). Here \(T=X_t\), and we shall accordingly write \(C_t\), \(E_t\) in place of \(C\), \(E\) above. One knows that \(C_t\) is projective over \(k(t)\), so (EGA IV 2.8.5) there is a projective flat scheme \(C\) over \(S\) whose generic fiber is \(C_t\). Using Abhyankar's resolution of singularities of arithmetic surfaces [6], one sees that one can even choose \(C\) to be a regular scheme. (N.B. Flatness is preserved under resolution, by EGA IV 14.3.8.) Since \(C\) is proper over \(S\), the points \(x_{0t}\) and \(x_{1t}\) of \(C_t\) rational over \(k(t)\) extend to sections \(x_0\), \(x_1\) of \(C\) over \(S\), and, since \(C\) is regular, it is well known that \(C\) is smooth over \(S\) at the points of \(x_0(S)\), \(x_1(S)\) (EGA IV, 17.12.1 c) \(\Rightarrow\) a) and 19.1.1). As a 